\begin{definition}[Primitive Dirichlet character] \label{definition:primitive_dirichlet_character}
\TODO{AI generated, verify}
Let $k \in \mathbb{Z}_{\geq 1}$, and let $\chi$ be a \CrefAndHyperrefIfExist{definition:dirichlet_character_modulo_a_nonnegative_integer}{Dirichlet character modulo $k$}. The character $\chi$ is called \hldef{primitive} if there does not exist any proper divisor $d$ of $k$ and any \CrefAndHyperrefIfExist{definition:dirichlet_character_modulo_a_nonnegative_integer}{Dirichlet character} $\psi$ modulo $d$ such that
$$\chi(n) = \psi(n) \quad \text{for all } n \in \mathbb{Z} \text{ with } \gcd(n,k) = 1.$$
If such a divisor $d$ and character $\psi$ exist, then $\chi$ is called \hldef{imprimitive}, and we say that $\chi$ is \hldef{induced by $\psi$}. The smallest positive integer $k^*$ for which there exists a primitive Dirichlet character modulo $k^*$ inducing $\chi$ is called the \hldef{conductor of $\chi$}\CrefIfExists{definition:conductor_of_a_dirichlet_character}.
\end{definition}
